\section{Binary Search to Find X in a Sorted Array}
\label{sec:bs-find-x}

\problemstatement{We are given an array $\textit{arr}$ of $n$ elements,
sorted in increasing order, and a target value $X$. We must determine the
index of $X$ in the array if it is present, and return $-1$ if it is not.
We must do this faster than the obvious $O(n)$ linear scan.}

\begin{patternbox}
\emph{``Sorted array''} appearing anywhere in a problem statement,
combined with a request to \emph{find/locate} a specific value or the
boundary of some property, is the single most reliable trigger for binary
search in this entire chapter. The reason a sorted array enables binary
search is precisely the monotonicity property from
Section~\ref{sec:bs-intro}: define $P(i) = (\textit{arr}[i] \ge X)$; because
the array is sorted, $P$ is \texttt{false} for every index before $X$'s
position and \texttt{true} from $X$'s position onward (or, if $X$ is
absent, from the position where elements first exceed $X$ onward) -- a
single crossover, which is exactly what binary search is built to find.
\end{patternbox}

\subsection*{Understanding the problem carefully}

If the array were unsorted, knowing that $\textit{arr}[\textit{mid}] \ne X$
would tell us nothing about where $X$ might be -- it could be anywhere
among the remaining elements, so we would have no choice but to check them
all. Sortedness changes this completely: if
$\textit{arr}[\textit{mid}] < X$, then because every element to the left of
$\textit{mid}$ is even smaller (the array is sorted in increasing order),
\emph{none} of them can equal $X$ either, so the entire left half can be
discarded in one step. Symmetrically, if $\textit{arr}[\textit{mid}] > X$,
the entire right half can be discarded. Only when
$\textit{arr}[\textit{mid}] = X$ do we stop, having found the answer.

\begin{ideabox}[Candidate Space and Predicate]
The candidate space is the set of indices $\{0, 1, \dots, n-1\}$. We
maintain a search window $[\textit{lo}, \textit{hi}]$, initially the whole
array, and repeatedly look at the midpoint, using the comparison with $X$ to
discard exactly the half of the window that provably cannot contain $X$.
\end{ideabox}

\subsection*{Why halving the search space each time is correct}

Each comparison at $\textit{mid}$ produces one of exactly three outcomes,
and in every outcome we either find the answer immediately or can prove
that an entire half of the remaining window is irrelevant. Since no element
is ever incorrectly discarded (the sortedness guarantee makes every
discard provably safe) and the window strictly shrinks by roughly half on
every iteration, the process must terminate, either by finding $X$ or by
the window becoming empty ($\textit{lo} > \textit{hi}$), at which point we
can safely conclude $X$ is not present.

\subsection*{Algorithm}

\begin{enumerate}
  \item Initialize $\textit{lo} \gets 0$, $\textit{hi} \gets n-1$.
  \item While $\textit{lo} \le \textit{hi}$:
  \begin{enumerate}
    \item $\textit{mid} \gets \textit{lo} + (\textit{hi}-\textit{lo})/2$.
    \item If $\textit{arr}[\textit{mid}] = X$: return $\textit{mid}$.
    \item If $\textit{arr}[\textit{mid}] < X$: $\textit{lo} \gets
          \textit{mid}+1$.
    \item Otherwise: $\textit{hi} \gets \textit{mid}-1$.
  \end{enumerate}
  \item If the loop ends without returning, return $-1$.
\end{enumerate}

\begin{algorithm}[H]
\caption{Binary Search}
\begin{algorithmic}[1]
\Procedure{BinarySearch}{$\textit{arr}, X$}
  \State $\textit{lo} \gets 0$, $\textit{hi} \gets n-1$
  \While{$\textit{lo} \le \textit{hi}$}
    \State $\textit{mid} \gets \textit{lo} + (\textit{hi}-\textit{lo})/2$
    \If{$\textit{arr}[\textit{mid}] = X$}
      \State \Return \textit{mid}
    \ElsIf{$\textit{arr}[\textit{mid}] < X$}
      \State $\textit{lo} \gets \textit{mid} + 1$
    \Else
      \State $\textit{hi} \gets \textit{mid} - 1$
    \EndIf
  \EndWhile
  \State \Return $-1$
\EndProcedure
\end{algorithmic}
\end{algorithm}

\subsection*{Dry run}

\dryrun{Take $\textit{arr} = [-1, 0, 3, 5, 9, 12]$, $X = 9$.}

\begin{itemize}
  \item $\textit{lo}=0,\textit{hi}=5$: $\textit{mid}=2$,
        $\textit{arr}[2]=3 < 9$. $\textit{lo}=3$.
  \item $\textit{lo}=3,\textit{hi}=5$: $\textit{mid}=4$,
        $\textit{arr}[4]=9 = 9$. Return $4$.
\end{itemize}

The answer is index $4$.

Now take $X = 2$ on the same array.

\begin{itemize}
  \item $\textit{lo}=0,\textit{hi}=5$: $\textit{mid}=2$,
        $\textit{arr}[2]=3 > 2$. $\textit{hi}=1$.
  \item $\textit{lo}=0,\textit{hi}=1$: $\textit{mid}=0$,
        $\textit{arr}[0]=-1 < 2$. $\textit{lo}=1$.
  \item $\textit{lo}=1,\textit{hi}=1$: $\textit{mid}=1$,
        $\textit{arr}[1]=0 < 2$. $\textit{lo}=2$.
  \item $\textit{lo}=2 > \textit{hi}=1$: loop ends. Return $-1$.
\end{itemize}

The answer is $-1$: $2$ is not present in the array.

\subsection*{C++ implementation (iterative and recursive)}

\begin{lstlisting}
#include <bits/stdc++.h>
using namespace std;

// Iterative version
int binarySearch(vector<int>& arr, int X) {
    int lo = 0, hi = arr.size() - 1;

    while (lo <= hi) {
        int mid = lo + (hi - lo) / 2;
        if (arr[mid] == X) return mid;
        else if (arr[mid] < X) lo = mid + 1;
        else hi = mid - 1;
    }
    return -1;
}

// Recursive version
int binarySearchRecursive(vector<int>& arr, int lo, int hi, int X) {
    if (lo > hi) return -1;

    int mid = lo + (hi - lo) / 2;
    if (arr[mid] == X) return mid;
    else if (arr[mid] < X) return binarySearchRecursive(arr, mid + 1, hi, X);
    else return binarySearchRecursive(arr, lo, mid - 1, X);
}

int main() {
    vector<int> arr = {-1, 0, 3, 5, 9, 12};
    cout << "Iterative: " << binarySearch(arr, 9) << endl;
    cout << "Recursive: "
         << binarySearchRecursive(arr, 0, arr.size() - 1, 9) << endl;
    return 0;
}
\end{lstlisting}

\begin{complexitybox}
\textbf{Time Complexity.} Each iteration (or recursive call) discards at
least half of the remaining search window, so the window shrinks from size
$n$ to size $0$ in at most $\lceil \log_2 n \rceil + 1$ steps, giving
\[
  O(\log n).
\]
\textbf{Space Complexity.} The iterative version uses only the three
scalar variables \textit{lo}, \textit{hi}, \textit{mid}, giving $O(1)$
space. The recursive version uses $O(\log n)$ space for the call stack,
since the recursion depth equals the number of halvings performed.
\end{complexitybox}

\begin{takeawaybox}
Binary search's correctness rests entirely on being able to \emph{prove},
at each step, that an entire half of the remaining search space cannot
contain the answer. In plain array search this proof comes directly from
sortedness; every other binary search problem in this chapter is really
asking you to find or construct an equivalent monotonicity proof in a
different setting.
\end{takeawaybox}
